\chapterbackmatter{Weinberg Exercises}

\begin{enumerate}
  \WeinbergExercise{1} Consider an $SU(N)$ gauge theory with a scalar field in the defining representation of $SU(N)$. Calculate the beta-function for the gauge coupling to one-loop order, including the contribution of a scalar loop. (Recommendation: Use the background field gauge, with a constant background field.)

  \WeinbergExercise{2} Suppose that the beta-function $\beta(g)$ for a theory with positive coupling constant $g$ has a simple zero at $g=g_*$, where $\beta(g)\to a(g_*-g)$ with $a>0$. What is the asymptotic behavior of the correction to the leading term $\propto E^{-\gamma^{\mathcal O}(g_*)}$ in the factor $N_E^{\mathcal O-1}$ associated with the inclusion of an operator $\mathcal O$ in a vacuum expectation value.

  \WeinbergExercise{3} Show that in a theory with $\beta(g)=bg^2+b'g^3+b''g^4+\cdots$, it is possible by a redefinition of the coupling constant to make the coefficient $b''$ anything we want.

  \WeinbergExercise{4} Calculate the effective electric charge that should be used in studying processes at energy 100 GeV, taking account of all known charged quarks and leptons with masses below 100 GeV.

  \WeinbergExercise{5} Calculate the asymptotic behavior for large four-momentum of the electron propagator in quantum electrodynamics. (You may use the one-loop value of $Z_2$ calculated elsewhere, for instance in Section 11.4.)

  \WeinbergExercise{6} Calculate the anomalous exponent $\nu$ to first order in the expansion in $\epsilon=4-d$ for an $O(N)$-invariant theory of scalar fields $\phi_n(x)$ with $n=1,\ldots,N$ belonging to the vector representation of $O(N)$, and an interaction $\tfrac14 g(\sum_n\phi_n^2)^2$.
\end{enumerate}
